% ==============================================================================
% SLIDES - SEMANA 05: GESTÃO DE SERVIÇOS, RACI, POLÍTICAS & ENTREGA PARCIAL 1
% ==============================================================================

% ------------------------------------------------------------------------------
% 1. SLIDE DE CAPA
% ------------------------------------------------------------------------------
\senaicover
  {Serviços, Estrutura \& Políticas}
  {Comitê Estratégico, Matriz RACI, Políticas de TI e Entrega Parcial 1}
  {Unidade Curricular: Governança de TI}
  {Prof. André Souza}

% ------------------------------------------------------------------------------
% 2. VISÃO GERAL / AGENDA
% ------------------------------------------------------------------------------
\begin{senaiframe}{Visão Geral \& Agenda}{Os 4 Eixos da Aula}

  \begin{columns}[t]
    \begin{column}{0.48\textwidth}
      \begin{tcolorbox}[senaibluecard, title={1. Conceito de Serviço}]
        \begin{itemize}
          \item Entrega de valor sem custos/riscos diretos.
          \item O Comitê Estratégico de TI (Steering).
        \end{itemize}
      \end{tcolorbox}
      
      \vspace{4pt}
      
      \begin{tcolorbox}[senaiambercard, title={3. Políticas Fundamentais}]
        \begin{itemize}
          \item Política de Segurança da Informação (PSI).
          \item Uso Aceitável de Ativos e Aquisições de TI.
        \end{itemize}
      \end{tcolorbox}
    \end{column}

    \begin{column}{0.48\textwidth}
      \begin{tcolorbox}[senairedcard, title={2. A Matriz RACI}]
        \begin{itemize}
          \item Responsible, Accountable, Consulted, Informed.
          \item Eliminação da dubiedade operacional.
        \end{itemize}
      \end{tcolorbox}
      
      \vspace{4pt}
      
      \begin{tcolorbox}[senaigreencard, title={4. Entrega Parcial 1}]
        \begin{itemize}
          \item Consolidação dos Capítulos 1 a 4.
          \item Submissão formal do primeiro marco do PDGTI.
        \end{itemize}
      \end{tcolorbox}
    \end{column}
  \end{columns}

\end{senaiframe}

% ------------------------------------------------------------------------------
% 3. MATRIZ RACI NOS PROCESSOS DE TI
% ------------------------------------------------------------------------------
\begin{senaiframe}{Módulo 1 • Estrutura de Governança}{Regra de Ouro da Matriz RACI}

  \begin{columns}[t]
    \begin{column}{0.48\textwidth}
      \begin{tcolorbox}[senaibluecard, title={Os 4 Papéis da Matriz RACI}]
        \begin{itemize}
          \item \textbf{R (Responsible):} Quem executa a tarefa na prática.
          \item \textbf{A (Accountable):} O único aprovador final (quem responde pelo sucesso/falha).
          \item \textbf{C (Consulted):} Especialista consultado antes de decidir.
          \item \textbf{I (Informed):} Pessoas avisadas após a execução.
        \end{itemize}
      \end{tcolorbox}
    \end{column}

    \begin{column}{0.48\textwidth}
      \begin{tcolorbox}[senairedcard, title={Regra de Ouro (Apenas 1 'A')}]
        \begin{itemize}
          \item Para cada atividade de TI, deve existir \textbf{exatamente um 'A'}.
          \item Se houver dois 'A's: Conflito de autoridade.
          \item Se houver zero 'A's: Ninguém responde pelas falhas.
        \end{itemize}
      \end{tcolorbox}
    \end{column}
  \end{columns}

\end{senaiframe}

% ------------------------------------------------------------------------------
% 4. POLÍTICAS FUNDAMENTAIS DE TI
% ------------------------------------------------------------------------------
\begin{senaiframe}{Módulo 2 • Diretrizes}{As 3 Políticas Fundamentais de TI}

  \begin{columns}[t]
    \begin{column}{0.48\textwidth}
      \begin{tcolorbox}[senaibluecard, title={Segurança \& Uso Aceitável}]
        \begin{itemize}
          \item \textbf{Política de Segurança (PSI):} Senhas, criptografia, MFA, acessos e resposta a incidentes.
          \item \textbf{Uso Aceitável:} Regras para notebooks, e-mail corporativo, internet e BYOD.
        \end{itemize}
      \end{tcolorbox}
    \end{column}

    \begin{column}{0.48\textwidth}
      \begin{tcolorbox}[senaiambercard, title={Aquisições \& Fornecedores}]
        \begin{itemize}
          \item \textbf{Política de Aquisição de TI:} Processo formal de cotação, homologação de softwares e SaaS.
          \item \textbf{Gestão de Terceiros:} Contratos de SLA, NDAs e gestão de riscos em fornecedores.
        \end{itemize}
      \end{tcolorbox}
    \end{column}
  \end{columns}

\end{senaiframe}

% ------------------------------------------------------------------------------
% 5. REFERÊNCIAS BIBLIOGRÁFICAS
% ------------------------------------------------------------------------------
\begin{senaiframe}{Referências Bibliográficas}{Fontes \& Leitura Recomendada}

  \begin{tcolorbox}[senaicard, title={Obras de Referência}]
    \begin{itemize}
      \item \textbf{FERNANDES, Aguinaldo Aragon; ABREU, Vladimir Ferraz.} *Implantando a governança de TI: da estratégia à gestão de processos e serviços.* 4. ed. Rio de Janeiro: Brasport, 2014.
      \item \textbf{FREITAS, Marcos André dos Santos.} *Fundamentos do gerenciamento de serviços de TI: preparatório para a certificação ITIL Foundation Edição 2011.* 2. ed. Rio de Janeiro: Brasport, 2018.
    \end{itemize}
  \end{tcolorbox}

\end{senaiframe}

% ------------------------------------------------------------------------------
% 6. APLICAÇÃO NO PDGTI
% ------------------------------------------------------------------------------
\begin{senaiframe}{Aplicação no PDGTI}{Orientações para a ENTREGA PARCIAL 1}

  \vspace{0.2cm}
  \begin{tcolorbox}[senaigreencard, title={Entrega Parcial 1 (Capítulos 1 a 4) – Data: 04/Set}]
    \begin{itemize}
      \item \textbf{Capítulo 1:} Caracterização da Organização e Estrutura Atual de TI.
      \item \textbf{Capítulo 2:} Diagnóstico da Situação Atual de TI (SWOT, dores e maturidade).
      \item \textbf{Capítulo 3:} Objetivos Estratégicos \& Mapa Estratégico (IT BSC).
      \item \textbf{Capítulo 4:} Estrutura de Governança, Comitês, Matriz RACI e Minuta de Políticas.
    \end{itemize}
  \end{tcolorbox}

\end{senaiframe}
